\begin{comment}

\begin{abstract}
Cloud\textendash{}Fog\textendash{}Edge infrastructures increasingly need to deploy and observe application logic across cloud services, gateways, and embedded Internet of Things (IoT) devices, but conventional cloud-native orchestration assumes Linux-class nodes with container runtimes, resident agents, and sufficiently rich networking. This assumption leaves a gap for constrained endpoints that must remain lightweight while still requiring secure attachment, life-cycle accountability, and reproducible evidence for sustainable research infrastructures. This work addresses the problem of expressing Kubernetes deployment intent for devices that cannot operate as Kubernetes nodes and translating that intent into verifiable embedded execution. We propose a Kubernetes-native framework that models devices, applications, and gateways through Custom Resource Definitions (CRDs) and uses a gateway-mediated control path to bind device identity, supervise liveness, dispatch WebAssembly (WASM) workloads, and report execution evidence back to the control plane. The gateway acts as the semantic boundary between declarative cloud state and constrained device protocols, preserving device simplicity while maintaining status ownership in Kubernetes. The framework is evaluated on a K3s-based test bed with emulated embedded targets and a complete enrollment-to-deployment workflow. % Removed: Across 100 independent trials, enrollment, heartbeat supervision, deployment, and transactional completion all succeeded (Wilson 95\% lower bound 96.3\%). Mean end-to-end latency is 1185.8~ms (95\% CI 1152.8--1218.7~ms). Measured southbound control traffic is 1728~B/h per device at a 25~s heartbeat period;
{\color{red}Across 100 repeated enrollment-to-deployment trials started from a known logical state, enrollment, heartbeat observability, deployment, and transactional completion all succeeded (Wilson 95\% lower bound 96.3\%). For the minimal 33~B WASM test module, mean end-to-end orchestration latency is 1286.0~ms (95\% CI 1226.4--1345.6~ms). Steady-state application-layer heartbeat traffic is 1728~B/h per device at a 25~s heartbeat period, excluding TLS and lower-layer overhead;} the gateway pod averages under 5~MiB RAM and 3~millicores CPU with one TLS-connected device, while {\color{red}the evaluated device firmware occupies 397~KiB of flash and approximately 264~KiB of static RAM, and the endpoint carries no kubelet or container runtime}. For sustainable digital research infrastructures, the framework therefore provides a reproducible, measurement-ready substrate for energy-aware placement, instrumentation, and green experimentation policies. % {\color{red}without claiming direct energy or carbon reduction in this evaluation}.
\end{abstract}
\end{comment}

\begin{abstract}
Cloud--Fog--Edge infrastructures increasingly need to deploy and observe application logic across cloud services, gateways, and embedded IoT devices, but conventional cloud-native orchestration assumes Linux-class nodes with container runtimes, resident agents, and rich networking. This leaves a gap for constrained endpoints that must remain lightweight while requiring secure attachment, life-cycle accountability, and reproducible evidence for sustainable research infrastructures. This work addresses the problem of expressing Kubernetes deployment intent for devices that cannot operate as Kubernetes nodes, and translating that intent into verifiable embedded execution. We propose a Kubernetes-native framework that models devices, applications, and gateways through Custom Resource Definitions (CRDs) and uses a gateway-mediated control path to bind device identity, supervise liveness, dispatch WebAssembly (WASM) workloads, and report execution evidence to the control plane. The gateway acts as the semantic boundary between declarative cloud state and constrained device protocols, preserving device simplicity while maintaining status ownership in Kubernetes. The framework is evaluated on a K3s-based test bed with emulated embedded targets across a complete enrollment-to-deployment workflow. Across 100 repeated trials, enrollment, heartbeat supervision, deployment, and transactional completion all succeeded. %(Wilson 95\% lower bound 96.3\%), with mean end-to-end latency of 1185.8~ms (95\% CI 1152.8--1218.7~ms). 
Heartbeat traffic stays under 2~KB/h per device, gateway pod footprint remains in the single-digit MiB and single-digit millicores, and the evaluated device firmware fits within 400~KiB of flash with no kubelet or container runtime. For sustainable research infrastructures, the framework provides a reproducible, measurement-ready substrate for energy-aware placement, instrumentation, and green experimentation policies.
\end{abstract}

\begin{keywords}
Cloud\textendash{}Fog\textendash{}Edge continuum \sep{} digital research infrastructures \sep{} sustainable orchestration \sep{} green experimentation \sep{} Internet of Things \sep{} WebAssembly \sep{} embedded systems \sep{} Kubernetes
\end{keywords}
